\section{Related Work}
\label{sec:related_work}

\subsection{Fine-Grained Animal Categorization and Datasets}
Fine-grained visual categorization (FGVC) focuses on differentiating subordinate categories characterized by high intra-class variance—caused by changes in pose, scale, and background clutter—and low inter-class variance, where discriminative traits are confined to localized regions. Standard benchmarks such as the Oxford-IIIT Pet dataset \cite{parkhi2012catsdogs} (37 cat and dog breeds) and Stanford Dogs \cite{khosla2011novel} (120 dog breeds, derived from ImageNet \cite{deng2009imagenet}) provide ground-truth bounding boxes and breed annotations for single-subject scenes. 

However, real-world deployment presents challenges absent from standardized benchmarks: co-occurring animals, background noise, and severe class imbalances. While recent wildlife models such as SpeciesNet \cite{speciesnet2025} and MegaDetector \cite{megadetector2026} tackle species localization in camera-trap settings, they are not designed for localized pet breed recognition under explicit operational constraints. Furthermore, relying on uncurated web scraping introduces severe label noise (e.g., commercial product names colliding with breed nomenclature). To overcome these limitations, we assemble a curated $35,423$-image multisource dataset structured specifically around the target 20-class taxonomy and an explicit reject class.

\subsection{Efficient Object Detection and Cascade Pipelines}
Single-stage object detectors, exemplified by the YOLO family, unify bounding box regression and class probability estimation into a single forward pass. \texttt{YOLOv6} \cite{li2022yolov6} introduces hardware-aware neural architecture design and efficient decoupled heads, optimizing execution speed for real-time industrial applications.

Rather than retraining an end-to-end detector across all fine-grained breeds—which requires dense, multi-class bounding-box annotations—our framework utilizes an off-the-shelf, COCO-pretrained \texttt{YOLOv6} \cite{lin2014coco} model strictly as a spatial localization and filtering stage. By isolating the spatially largest target bounding box and applying zero-out spatial masking to secondary detections, our cascade mitigates background context leakage and reduces the downstream classification task to standardized single-subject crops.

\subsection{Efficient Backbones and Spatial Regularization}
Convolutional backbones designed for mobile and real-time deployment must balance parameter footprint with representational capacity. \texttt{EfficientNetV2} \cite{tan2021efficientnetv2} addresses this by combining Neural Architecture Search (NAS) with fused Mobile Inverted Bottleneck (Fused-MBConv) layers, significantly accelerating training throughput compared to original scaling laws \cite{tan2019efficientnet}.

When training deep networks from scratch on specialized FGVC tasks, models are particularly susceptible to rapid validation loss divergence and memorization \cite{prechelt1998early}. To enforce spatial invariance, data augmentation strategies such as CutMix \cite{yun2019cutmix} and Random Erasing \cite{zhong2020random} blend localized image patches during training. In this work, we deploy \texttt{EfficientNetV2-S} backbones and demonstrate that spatial patch-mixing (\texttt{CropMix}) serves as the primary regularizer, stabilizing from-scratch convergence and closing the generalization gap.

\subsection{Explainable AI and Background Shortcut Learning}
Deep Convolutional Neural Networks (CNNs) frequently exploit unintended visual shortcuts \cite{geirhos2020shortcut}, relying on background co-occurrences (e.g., grass, indoor furniture) rather than anatomical features—a phenomenon known as the "Clever Hans" effect \cite{lapuschkin2019unmasking}. 

Gradient-weighted Class Activation Mapping (Grad-CAM) \cite{selvaraju2017grad} generates coarse visual localization maps by computing the gradients of a target class score with respect to the final convolutional feature maps. In our framework, we employ Grad-CAM not merely as a visualization tool, but as a diagnostic mechanism to empirically verify that spatial zero-masking successfully forces the downstream classification heads to focus on facial and morphological breed characteristics rather than background distractors.